% Tabla de resultados Avance 5 — pesos re-entrenados en SCARED (Dr. Espinosa Loera)
% Equipo 52 — Proyecto Integrador TC5035.10
% Requiere: \usepackage{booktabs, multirow}
\begin{table}[t]
\centering
\caption{Avance 5: Resultados de rendimiento y comparación de métodos de mejora (Enhancement) frente al baseline, usando pesos re-entrenados en SCARED.}
\label{tab:avance5}
\setlength{\tabcolsep}{4pt}
\renewcommand{\arraystretch}{1.05}
\small
\begin{tabular}{llcccccccrr}
\toprule
\textbf{Modelo} & \textbf{Enhancement} & $\epsilon_{\text{AbsRel}}\!\downarrow$ & $\epsilon_{\text{SqRel}}\!\downarrow$ & $\epsilon_{\text{RMSE}}\!\downarrow$ & $\epsilon_{\text{RMSELog}}\!\downarrow$ & $\delta\!<\!1.25\!\uparrow$ & $\delta\!<\!1.25^2\!\uparrow$ & $\delta\!<\!1.25^3\!\uparrow$ & IT (ms) & $\Delta$AR\% \\
\midrule
\multirow{4}{*}{Monodepth2}
 & none       & 0.0568 & 0.6821 &  4.771 & 0.0777 & 0.9852 & 0.9915 & 0.9939 &  27.4 & baseline \\
 & retinex    & 0.1114 & 1.4798 &  8.236 & 0.1489 & 0.8656 & 0.9742 & 0.9897 & 413.0 & +96.1\% \\
 & endolmspec & 0.0662 & 0.8576 &  5.324 & 0.0876 & 0.9787 & 0.9912 & 0.9929 & 150.0 & +16.5\% \\
 & iat        & 0.0581 & 0.6269 &  4.796 & 0.0791 & 0.9832 & 0.9913 & 0.9948 & 285.3 & +2.3\% \\
\midrule
\multirow{4}{*}{MonoViT}
 & none       & 0.1328 & 1.9229 &  9.966 & 0.1752 & 0.8476 & 0.9685 & 0.9868 &  51.4 & baseline \\
 & retinex    & 0.1339 & 1.9151 &  9.809 & 0.1755 & 0.8397 & 0.9535 & 0.9887 & 437.5 & +0.8\% \\
 & endolmspec & 0.1425 & 2.1251 & 10.541 & 0.1841 & 0.8173 & 0.9666 & 0.9864 & 173.9 & +7.3\% \\
 & iat        & 0.1353 & 1.9456 &  9.930 & 0.1758 & 0.8464 & 0.9643 & 0.9874 & 309.4 & +1.9\% \\
\midrule
\multirow{4}{*}{EndoSfMLearner}
 & none       & 0.1706 & 3.1434 & 12.754 & 0.2074 & 0.7750 & 0.9387 & 0.9884 &   6.4 & baseline \\
 & retinex    & 0.1710 & 3.1565 & 12.777 & 0.2077 & 0.7743 & 0.9386 & 0.9884 & 392.2 & +0.2\% \\
 & endolmspec & 0.1707 & 3.1466 & 12.761 & 0.2074 & 0.7747 & 0.9387 & 0.9884 & 129.1 & +0.1\% \\
 & iat        & 0.1707 & 3.1476 & 12.761 & 0.2075 & 0.7747 & 0.9387 & 0.9884 & 264.3 & +0.1\% \\
\midrule
\multirow{4}{*}{AF-SfMLearner}
 & none       & 0.1103 & 1.3733 &  8.970 & 0.1420 & 0.9008 & 0.9870 & 0.9966 &  27.2 & baseline \\
 & retinex    & 0.1326 & 1.8277 &  9.682 & 0.1786 & 0.8552 & 0.9600 & 0.9874 & 413.1 & +20.2\% \\
 & endolmspec & \underline{0.1055} & \underline{1.2379} & \underline{8.483} & \underline{0.1365} & \underline{0.9179} & 0.9869 & \underline{0.9968} & 150.2 & $-$4.4\% $\leftarrow$ \\
 & iat        & 0.1079 & 1.2953 &  8.678 & 0.1408 & 0.9026 & 0.9844 & 0.9967 & 285.3 & $-$2.2\% $\leftarrow$ \\
\midrule
\multirow{4}{*}{MonoIIT}
 & none       & \textbf{0.0373} & \textbf{0.2642} & \textbf{3.487} & \textbf{0.0572} & \textbf{0.9886} & 0.9958 & 0.9995 &  72.7 & baseline \\
 & retinex    & 0.0988 & 1.1629 &  7.732 & 0.1353 & 0.8911 & 0.9822 & 0.9947 & 458.5 & +164.9\% \\
 & endolmspec & 0.0424 & 0.3540 &  3.700 & 0.0618 & 0.9817 & 0.9926 & 0.9985 & 195.5 & +13.7\% \\
 & iat        & 0.0411 & 0.3067 &  3.637 & 0.0600 & 0.9848 & 0.9941 & 0.9994 & 330.7 & +10.2\% \\
\midrule
\multirow{4}{*}{w19-MonoViT}
 & none       & 0.0459 & 0.2755 &  3.913 & 0.0637 & 0.9881 & \textbf{0.9986} & \textbf{0.9997} &  72.5 & baseline \\
 & retinex    & 0.1342 & 2.2202 &  9.920 & 0.1850 & 0.8114 & 0.9305 & 0.9740 & 458.4 & +192.4\% \\
 & endolmspec & 0.0597 & 0.3932 &  4.705 & 0.0786 & 0.9827 & 0.9984 & 0.9997 & 195.4 & +30.1\% \\
 & iat        & 0.0524 & 0.3212 &  4.108 & 0.0689 & 0.9880 & 0.9978 & 0.9996 & 330.6 & +14.2\% \\
\bottomrule
\end{tabular}

\vspace{2pt}
\footnotesize
\textbf{Nota:} $\delta < 1.25^k$ representa la fracción de píxeles con $\max(\text{pred}/\text{gt}, \text{gt}/\text{pred}) < 1.25^k$ (valores más altos indican mejor rendimiento). Un porcentaje negativo en $\Delta$AR\% indica una mejora (\emph{improvement}) respecto a la configuración sin procesar (\emph{none}). Los casos que logran esta mejora se señalan con el símbolo $\leftarrow$. Pesos re-entrenados en SCARED proporcionados por el Dr. Ricardo Espinosa Loera. \textbf{Negrita} = mejor valor global; \underline{subrayado} = mejor enhancement por modelo.
\end{table}
